\section{Materials and Methods}

\subsection{Data and Software Environment}
The analysis uses one apical four-chamber video and three short-axis videos (MV, PM, AP), together with Doppler measurements and heart rate recorded during the laboratory sessions. Processing was implemented in Python using a dedicated application (\texttt{main.py}, \texttt{heartvolume/gui/app.py}) and numerical routines in \texttt{heartvolume/core/calculations.py}.

\subsection{Handout-Driven Analysis Logic}
Following the Week 2\&3 instructions, the report compares three cardiac-function techniques: (i) Week 2 Simple Method for EDV/ESV, (ii) Week 2 Doppler method for flow-derived CO, and (iii) Week 3 Modified Simpson method for EDV/ESV. For volume-based methods (Simple and Simpson), SV, EF, and CO are computed from EDV and ESV. Doppler CO is computed independently from $VTI$, $HR$, and $D_{AO}$.

\subsection{Week 2: Simple Method (2-D Geometric Approximation)}
The first ventricular-volume approach is referred to as the \emph{Simple Method}. Assuming the ventricular cavity as a truncated axisymmetric ellipsoid, volume is estimated by
\begin{equation}
V = \frac{\pi}{6} D_{\mathrm{endo}}^2 L_{\mathrm{endo}}.
\end{equation}
The same formula was applied at end-diastole and end-systole to obtain $EDV_{\mathrm{simple}}$ and $ESV_{\mathrm{simple}}$, consistent with standard echocardiographic quantification practice \cite{schiller1989lv}.

\subsection{Week 2: Balloon Phantom Formula}
The balloon is modeled as an ellipsoid with major axis $L$ and orthogonal minor axes $D_1$ and $D_2$:
\begin{equation}
V_{\mathrm{balloon}} = \frac{4}{3}\pi\left(\frac{L}{2}\right)\left(\frac{D_1}{2}\right)\left(\frac{D_2}{2}\right) = \frac{\pi}{6} L D_1 D_2.
\end{equation}
In this report, four sets of measurements are provided. Each set is used to compute an estimated balloon volume, then absolute and relative errors versus the known injected volume.

\subsection{Week 2: Doppler-Derived Cardiac Output}
Cardiac output from aortic outflow was computed as
\begin{equation}
CO_{\mathrm{Doppler}} = VTI \times HR \times \frac{\pi D_{AO}^2}{4}.
\end{equation}
The aortic diameter was measured directly on the D\_AO image/video after calibration, and VTI and HR were entered from the acquisition session.

\subsection{Week 3: Modified Simpson Method Implemented in Python}
The Python workflow reproduces the Week 3 procedure with interactive frame selection and calibrated manual measurements:
\begin{itemize}
    \item frame extraction from the four synchronized views;
    \item automatic or manual scale calibration in cm/pixel;
    \item manual distance measurements on ED and ES frames;
    \item modified Simpson volume computation through \texttt{calculate\_volume\_simpson}.
\end{itemize}
For each cardiac phase, the implementation uses long-axis length $L$ and the three short-axis diameters at apex, papillary muscle, and mitral valve levels. Using areas $A_{\mathrm{AP}}$, $A_{\mathrm{PM}}$, and $A_{\mathrm{MV}}$,
\begin{equation}
V_{\mathrm{Simpson}} = \frac{L}{3}\times\frac{A_{\mathrm{AP}} + A_{\mathrm{MV}} + 4A_{\mathrm{PM}}}{3},
\end{equation}
with
\begin{equation}
A_i = \pi\left(\frac{D_i}{2}\right)^2, \quad i\in\{\mathrm{AP},\mathrm{PM},\mathrm{MV}\}.
\end{equation}
This provides $EDV_{\mathrm{simpson}}$ and $ESV_{\mathrm{simpson}}$.

\subsection{Derived Metrics and Comparison Framework}
For each volume-based method (Simple and Modified Simpson), stroke volume and ejection fraction were computed as
\begin{align}
SV &= EDV - ESV, \\
EF &= \frac{SV}{EDV} \times 100, \\
CO &= SV \times HR.
\end{align}
Thus, three cardiac output estimates are compared in this report:
\begin{itemize}
    \item $CO_{\mathrm{simple}} = SV_{\mathrm{simple}} \times HR$,
    \item $CO_{\mathrm{simpson}} = SV_{\mathrm{simpson}} \times HR$,
    \item $CO_{\mathrm{doppler}} = VTI \times HR \times (\pi D_{AO}^2/4)$.
\end{itemize}
The full comparison is therefore organized as follows:
\begin{itemize}
    \item compare ventricular volumes from Simple Method vs Modified Simpson,
    \item compare cardiac output from Simple vs Modified Simpson vs Doppler,
    \item compare balloon estimated volume vs true balloon volume (250 mL).
\end{itemize}
In the Python code, \texttt{build\_volume\_result} and \texttt{build\_doppler\_result} centralize these calculations and standardize units (mL and L/min) for consistent reporting across methods.

\subsection{Error Analysis Strategy}
Method comparison focused on sensitivity to manual annotations, frame selection, calibration quality, and geometric assumptions. Differences between methods were interpreted by considering expected physiological ranges and known acquisition-related sources of bias.
